\chapter{Justificación}\label{cap2}
Las partículas de polvo en el área de trabajo de un panel solar, y la frecuencia en la que se limpian los paneles solares puede llegar a incrementar la pérdida de generación de energía, lo que nos lleva a querer resolver la problemática que tienen los paneles solares que es el hecho de que su ciclo operativo se ve afectado en su funcionamiento óptimo por agentes externos como el polvo, heces de pájaros, nubes además de que se puede considerar que al limpiarlo se hace un efecto de sombra con el elemento con el que sea limpiado, aparte de tomar en cuenta el costo de limpieza manual de los módulos solares.\\

Aparte es difícil mantener a una persona al pendiente del cuidado y limpieza de un panel solar, se crean complicaciones y costos dados por que se tienen que hacer pagos a la persona por el tiempo requerido cuidando que el panel esté limpio y por cada panel que se requiera limpiar. costo de limpieza manual por un panel solar está generalmente dado en 0.12 dlrs [11].\\

Las inherentes limitaciones existentes en los materiales fotovoltaicos como lo es el polvo en las fachadas de los paneles solares, hace que se entorpezca el proceso de transmisividad de la luz, lo que afecta directamente la eficiencia del sistema y el desempeño en general.\\

Considerando esto, para que el sistema fotovoltaico pueda captar la mayor cantidad de energía solar y por lo tanto trabajar a su máxima capacidad, es necesario mantener la superficie de los paneles solares libre de elementos que puedan obstruir la captación de la luz solar.\\

Este proyecto se realiza con la finalidad de generar un mecanismo que pueda realizar la limpieza de conjuntos de paneles solares de manera periódica así logrando que la superficie de los paneles pueda estar libre de impedimentos que le impidan aprovechar la luz solar para convertirla en energía.



% \begin{figure}[H]
% 	\centering
% 	\includegraphics[scale=1]{upiiz}
% 	\caption{Este es un título muy largo para una Figura.}
% 	\label{fig2}
% \end{figure}


% \begin{equation}\label{key}
% \left [
% \begin{array}{c}
% \dot{x}\\
% \dot{y} \\
% \dot{\theta}
% \end{array}
% \right ]
% \end{equation}

% \begin{align}\label{key}
% \sum_{i}^{m}&=0\\
% \sum&=10
% \end{align}



% \begin{table}[H]	
% 	\centering
% 	\caption{Tabla demostrativa 2.}
% 	\begin{tabular}{|c|c|c|}
% 		\hline
% 		Celda 1,1&Celda 1,2&Celda 1,3\\\hline
% 		Celda 2,1&Celda 2,2&Celda 2,3\\\hline
% 	\end{tabular}
% \end{table}